\input{../tex-inputs/latex-leseansicht-vorspann}

\section[Arthur Schnitzler an Gustav Schwarzkopf, {[}4{]}. 3. 1900]{L04408 Arthur Schnitzler an Gustav Schwarzkopf, {[}4{]}. 3. 1900}
\nopagebreak\mylabel{L04408v}
\rehead{ }\normalsize\beginnumbering\briefempfaengerindex{Schwarzkopf, Gustav@\textsc{Schwarzkopf, Gustav}!zzzSchnitzler, Arthur@\emph{von Arthur Schnitzler}!1900-03-041@{{[}4{]}. 3. 1900}|(be}
\toendnotes[C]{\smallbreak\pagebreak[2]}
\correspDesc{Versand  durch Arthur Schnitzler am [4]. 3. 1900 in Edlach an der Rax
\newline{}Übermittlung  am 4. 3. 1900 in Edlach an der Rax
\newline{}Zustellung  am 5. 3. 1900 in Wien
\newline{}Erhalt  durch Gustav Schwarzkopf am [5]. 3. 1900 in Wien}\toendnotes[C]{\smallbreak}
\Standort{DLA, A:Schnitzler, HS.1985.1.1897.}
\physDesc{Bildpostkarte, 130 Zeichen
\newline{}Handschrift: Bleistift, deutsche Kurrent
\newline{}Versand: 1) Stempel: »\nobreak{}\oindex{Reichenau an der Rax@\textbf{Reichenau an der Rax}|pwk}Reichenau, 4/3 00\nobreak{}«.   2) Stempel: »\nobreak{}\oindex{I., Innere Stadt@\textbf{I., Innere Stadt}|pwk}Wien 1/1 1, 5. 3. 00, 8–9½V., Bestellt\nobreak{}«. }\pstart{}{\pb}\textsc{Herrn Gustav
                     Schwarzkopf}\pend{}\pstart{}Wien\oindex{Wien@\textbf{Wien}|pw}\pend{}\pstart{}\textsc{I. Tiefer Graben
                     23}\oindex{Wien@\textbf{Wien}!I., Innere Stadt@\textbf{I., Innere Stadt}!Tiefer Graben 23@\textbf{Tiefer Graben 23}|pw}\pend{}{\bigskip}
\pstart
           \centering{}\textcolor{gray}{\textbf{{\pb}Gruss aus Edlach\oindex{Edlach@\textbf{Edlach}|pw}}}\pend
           \vspace{1em}
\pstart
           \centering{}{\pb}\textcolor{gray}{\textbf{EDLACHER\oindex{Edlach@\textbf{Edlach}|pw}
                        WOHLTHÄTIGKEITSFEST}}\pend
           
\pstart
           \centering{}(Iſt gar  nicht wahr)\pend
           \vspace{0.5em}
\pstart
           Tiefſter Winter Aber ſehr ſchön und \textcolor{gray}{ſtill}\pend
           
\pstart
           Herzl. Grüße.{\\[\baselineskip]}Ihr \spacefill\mbox{Arth}.\pend
           \leftskip=0em{}\selectlanguage{ngerman}\endnumbering\briefempfaengerindex{Schwarzkopf, Gustav@\textsc{Schwarzkopf, Gustav}!zzzSchnitzler, Arthur@\emph{von Arthur Schnitzler}!1900-03-041@{{[}4{]}. 3. 1900}|)be}\mylabel{L04408h}
\begin{anhang}
\end{anhang}\newcommand{\dateiname}{L04408}\newcommand{\titel}{Arthur Schnitzler an Gustav Schwarzkopf, [4]. 3. 1900}\newcommand{\editorInnen}{Herausgegeben von Jahnke, SelmaMüller, Martin Anton}\input{../tex-inputs/latex-leseansicht-abspann}
